\subsection{Analogy with Programming Languages}
\label{sec:programming-analogy}
\slidesrc{Terminology introduction: 25--26}

The three levels have a natural counterpart in programs:

\begin{center}
\begin{tabular}{@{}l l@{}}
\toprule
\textbf{Level} & \textbf{A program is correct at this level if\ldots} \\
\midrule
\textcolor{nbred}{Syntax}      & there are no compiler errors \\
\textcolor{nbblue}{Semantics}  & there are no implementation bugs \\
\textcolor{nbgreen}{Pragmatics} & it implements the right algorithm (for the purpose) \\
\bottomrule
\end{tabular}
\end{center}

Three examples show that the levels are independent of each other:

\begin{itemize}
  \item \textbf{Different \textcolor{nbred}{syntax}, same \textcolor{nbblue}{semantics}}
        (both mean 5):
        \[ 2 + 3 \quad\Longleftrightarrow\quad 3 + 2 \]
  \item \textbf{Same \textcolor{nbred}{syntax}, different \textcolor{nbblue}{semantics}}
        (1 and 1.5): the expression \texttt{3 / 2} is \texttt{1} in Python~2.7 (integer
        division between integers) but \texttt{1.5} in Python~3 (true division):
        \[ \texttt{3 / 2}\ \text{(Python 2.7)} \quad\not\Longleftrightarrow\quad
           \texttt{3 / 2}\ \text{(Python 3)} \]
  \item \textbf{Good \textcolor{nbblue}{semantics}, bad \textcolor{nbgreen}{pragmatics}}:
        a correct, bug-free implementation of a deep neural network for estimating the
        probability of a coin flip. It works, but it is absurd for the purpose: counting
        heads and dividing by the number of flips does the job.
\end{itemize}

\begin{nbintuition}
Syntax is about the \emph{form}, semantics about the \emph{content}, pragmatics about
the \emph{use}. Each level can be right while the one above is wrong: a sentence (or a
program) can be well formed and meaningless, or meaningful and pointless.
\end{nbintuition}
